% 48-23(48쪽) — 곡선 y=x^3-5x+6 과 점 P(1,2) 에서의 접선. 접선이 곡선과 다시 만나는 점 Q, x축과 만나는 점 R. 스캔 화질이 낮아 TikZ 로 다시 그림.
% 원문에 있는 것만: 눈금 숫자 1 (x) · 2, 6 (y) · 라벨 O, x, y, P, Q, R, y=x^3-5x+6 · P 의 좌표 안내 파선.
% Q, R 의 좌표는 원문에 적혀 있지 않으므로 그림에도 적지 않는다.
\begin{tikzpicture}[x=0.38cm, y=0.20cm, line width=0.5pt]
  \draw[->, >=stealth, line width=0.7pt] (-3.6,0) -- (4.0,0) node[below=1pt] {$x$};
  \draw[->, >=stealth, line width=0.7pt] (0,-2.3) -- (0,11.9) node[left=1pt] {$y$};
  \node[below left=-1pt and -2pt, font=\small] at (0,0) {$\pt{O}$};
  % 좌표 안내 파선(원문)
  \draw[dashed, line width=0.35pt] (0,2) -- (1,2) -- (1,0);
  % 점 P 에서의 접선
  \draw[line width=0.5pt] (-3.45,10.9) -- (2.35,-0.7);
  % 곡선 y=x^3-5x+6
  \draw[line width=0.6pt, smooth] plot coordinates {
    (-2.80,-1.952) (-2.70,-0.183) (-2.60,1.424) (-2.40,4.176) (-2.20,6.352)
    (-2.00,8.000) (-1.80,9.168) (-1.60,9.904) (-1.40,10.256) (-1.29,10.303)
    (-1.20,10.272) (-1.00,10.000) (-0.80,9.488) (-0.60,8.784) (-0.40,7.936)
    (-0.20,6.992) (0.00,6.000) (0.20,5.008) (0.40,4.064) (0.60,3.216)
    (0.80,2.512) (1.00,2.000) (1.29,1.697) (1.50,1.875) (1.70,2.413)
    (1.90,3.359) (2.10,4.761) (2.30,6.667) (2.45,8.456) (2.60,10.576)};
  % 눈금 라벨(원문 자리 그대로: 2, 6 은 y축 왼쪽)
  \node[left=2pt, font=\small] at (0,2) {$2$};
  \node[left=2pt, font=\small] at (0,6) {$6$};
  \node[below=2pt, font=\small] at (1,0) {$1$};
  \node[above right=-1pt and -1pt, font=\small] at (1,2) {$\pt{P}$};
  \node[below left=-1pt and 0pt, font=\small] at (-2,8) {$\pt{Q}$};
  \node[above right=-1pt and 0pt, font=\small] at (2,0) {$\pt{R}$};
  \node[right=1pt, font=\small] at (0.2,11.2) {$y=x^3-5x+6$};
\end{tikzpicture}
